% Generated draft: requires mathematical and source review.
\documentclass[11pt]{article}
\usepackage[margin=25mm]{geometry}
\usepackage{fontspec}
\setmainfont{Latin Modern Roman}
\usepackage{amsmath,amssymb,mathtools,unicode-math}
\setmathfont{Latin Modern Math}
\usepackage{enumitem,needspace,xurl,hyperref}
\setlength{\parindent}{0pt}
\setlength{\parskip}{5pt}
\setlength{\emergencystretch}{3em}
\newcommand{\N}{\mathbb{N}}
\newcommand{\Z}{\mathbb{Z}}
\newcommand{\Q}{\mathbb{Q}}
\newcommand{\R}{\mathbb{R}}
\newcommand{\C}{\mathbb{C}}
\newcommand{\F}{\mathbb{F}}
\newcommand{\A}{\mathbb{A}}
\newcommand{\PP}{\mathbb{P}}
\newcommand{\E}{\mathbb{E}}
\newcommand{\eps}{\varepsilon}
\newcommand{\dd}{\,\mathrm{d}}
\begin{document}
\section*{UnsolvedMath ID: 30001243}
\textbf{Source problem code:} catalogue entry
\subsection{Definitions and mathematical statement}
Let \(X\) be a normal projective variety. The supplied record does not specify its ground field or characteristic. Formally, set
\[
A(X)=\bigl\{a(R)\in\Z:\ R\text{ is a quasi-Gorenstein graded ring and }\operatorname{Proj}R\cong X\bigr\}.
\]
The precise class of admissible graded rings is not specified in the supplied record; in particular, the grading, finiteness assumptions, and degree-zero ring require confirmation from the source.

In the usual graded canonical-module formulation, a ring \(R\) is quasi-Gorenstein when its graded canonical module satisfies
\[
\omega_R\cong R(a)
\]
for an integer \(a\). With the shift convention \(R(a)_n=R_{n+a}\), this integer is its \(a\)-invariant, denoted \(a(R)\). This explains the notation under that convention; the hypotheses ensuring existence of \(\omega_R\) and the intended convention must be checked against the source.

Let \(K_X\) denote a canonical divisor. For a \(\Q\)-Cartier divisor \(D\), bigness means that its Iitaka dimension equals \(\dim X\). Under the usual Cartier interpretation, a Fano variety has ample Cartier anticanonical divisor \(-K_X\), and its index is the largest positive integer \(r\) such that \(-K_X\sim rH\) for a Cartier divisor \(H\). The supplied record does not specify whether these Cartier hypotheses or other conventions for singular varieties are intended.

\begin{enumerate}
\item Classify the possible sets \(A(X)\) as \(X\) ranges over the intended normal projective varieties. More specifically, the original statement asks whether every \(A(X)\) equals one of the six sets listed in Example~5 of the cited report. Those six sets are absent from the supplied record, so this specific classification target cannot be reproduced explicitly.
\item Whenever \(-K_X\) is big, must there exist an admissible quasi-Gorenstein graded ring \(R\) with \(\operatorname{Proj}R\cong X\) and \(a(R)<0\)? Equivalently, must
\[
A(X)\cap\Z_{<0}\ne\varnothing?
\]
\item For every Fano variety \(X\) of index \(1\) in the intended class, is
\[
A(X)=\{2m+1:m\in\Z,\ m\geq0\}?
\]
\end{enumerate}
These are questions, not asserted conclusions. Under the usual Fano convention, \(-K_X\) is big, so the proposed affirmative answers to the second and third questions would be incompatible for the same index-one Fano variety. Both questions are retained as supplied, pending source review.
\subsection{Short English statement}
For each normal projective variety \(X\), collect the \(a\)-invariants of quasi-Gorenstein graded rings whose projective schemes are isomorphic to \(X\). Classify these sets, test whether a big anticanonical divisor forces a negative member, and ask whether an index-one Fano variety gives exactly the positive odd integers. The six proposed classification sets and several foundational conventions are missing from the supplied record.
\subsection{Sources}
\begin{itemize}
\item Supplied local problem registry record: permanent UnsolvedMath ID \(30001243\), source problem code \texttt{catalogue entry}, title \emph{a-Invariants of Quasi-Gorenstein Projective Cones}. The record belongs to the Oberwolfach Reports Open Problems collection and gives proposed year 2009. No individual proposer is supplied.
\item \emph{Kommutative Algebra} (2009), Oberwolfach Reports, Report 22, pp.~1157--1224, DOI \texttt{10.4171/owr/2009/22}. Supplied report reference: \url{https://doi.org/10.4171/owr/2009/22}.
\item Supplied literature source labelled ``Oberwolfach Report 22, Kommutative Algebra'': \url{https://doi.org/10.4171/owr/2009/22}.
\item Supplied literature source labelled ``original A(X) questions'': \url{https://doi.org/10.4171/owr/2009/22}.
\item Literature assessment dated exactly \texttt{2026-08-21}: the supplied curation labels the questions open and reports many computed cases but no general classification or criterion of the requested kind. This is supplied dated curation, not an independently verified literature assessment.
\item The registry labels the statement \texttt{corrected\_verified} and states that its mathematical content was checked against the official report. That verification claim is supplied metadata; it has not been independently checked for this draft.
\end{itemize}
\end{document}
